% title:          Gold and silver coins
% area:           processes
% methods:
% difficulty:
% difficulty-ai:  easy
% status:         partial
% review:         none
% --- history: all optional, leave EMPTY if unknown ---
% origin:         tournament-of-towns
% origin-date:    2014-10-12
% origin-number:  Fall, Senior O-Level, Problem 5
% also-in:
% taken-from:
% references:     Tournament of Towns, 36th tournament, Fall 2014 (12 Oct 2014), grades 10-11 basic variant (Senior O-Level), Problem 5 [5 points], author E. V. Bakaev - https://www.turgor.ru/problems/36/index.php; English paper 'Senior O-Level Paper Fall 2014', Problem 5 - https://www.math.toronto.edu/oz/turgor/archives/TT2014F_SOproblems.pdf; official solutions (compiled by L. Mednikov, area-under-staircase argument) - http://ashap.info/Turniry/TG/36ob.pdf; problems.ru, Problem 64845 - https://problems.ru/view_problem_details_new.php?id=64845; Math.SE 969781 'A beautiful game of gold and silver coins' (posted 12 Oct 2014, answers incl. Sean Lo's area proof) - https://math.stackexchange.com/questions/969781/a-beautiful-game-of-gold-and-silver-coins. Club source: Sudakov & Milojević (ETH), not located online.
% history:        ai
\begin{problem}
There is a pile of silver coins on a table. John holds two pieces of paper and performs the
following process: at each step, he can add one gold coin to the table and write the current number of silver
coins on one piece of paper, or remove one silver coin from the table and write down the current number
of gold coins on the other piece of paper. This process runs until no more silver coins remain on the table.
Show that at the end of the process, the sums of the numbers on both pieces of paper are equal.
\end{problem}
\begin{solution}
\[
\substack{\textit{This solution needs to be written more formally}}
\]
See e.g. the solution proposed by Sean Lo at this \href{https://math.stackexchange.com/questions/969781/a-beautiful-game-of-gold-and-silver-coins}{Post}
\end{solution}
